\section{Conclusiones}

La propuesta se ha convertido en un estudio completo: dataset auditado, tres
modelos, protocolo sin fuga, nueve ejecuciones científicas, controles simples y
análisis de errores. La contribución defendible no es afirmar que un
autoencoder diagnostica patologías, sino medir bajo qué condiciones la hipótesis
de reconstrucción separa normalidad y no normalidad.

GANomaly es el mejor modelo neuronal: supera AE y VAE en las tres semillas y
alcanza AUROC $0{,}5636\pm0{,}0150$. AE y VAE no son detectores útiles en esta
configuración, pues quedan consistentemente por debajo de 0,5 aunque reconstruyen
bien las normales. Sin embargo, el control de gradiente obtiene 0,5981 y
evidencia que variaciones simples de textura o adquisición explican una parte
importante de la separación. No se sostiene que GANomaly haya aprendido una
señal clínica robusta ni que esté preparado para uso asistencial.

El resultado es útil y defendible: demuestra experimentalmente los límites de
la hipótesis de reconstrucción, cuantifica el beneficio y coste del entrenamiento
adversarial, y detecta un atajo que quedaría oculto reportando solo AUPRC. La
conclusión negativa respecto a utilidad clínica evita una afirmación
injustificada.

Como trabajo futuro se priorizan una ablación a 128$\times$128, validación
externa por pacientes y comparación con representaciones preentrenadas o
autoencoders enmascarados \cite{he2022mae,georgescu2023masked}. Solo después
tendría sentido estudiar localización o despliegue.
